\par
The main difference between machine learning and statistics is that in machine
learning one takes the data to predict new data, whilst in statistics this is
not done, it is done for missing data points within existing data.

\section{Dataset splitting}

\par
Train, validation and final test.

% TODO: Scikit example

\par
The very first thing we sholud do is calculate the \textit{out-of-sample}
performance of the \textit{baseline} model.

$$
	\hat{y} = \mathbb{E}[f(x)] = \overline{y}\quad \forall \hat{y}
$$

% TODO: Include plot

% TODO: Include baseline code and theory

$$
	RMSE = \sqrt{\overline{\epsilon^2}} =  \sqrt{\frac{1}{n}\sum(y_i - \hat{y_i})^2}
$$

\par
One can asses the model visually by plotting the predicted values against the
true values (\textit{true} vs \textit{predicted} plot). The ideal situation in
this case is for the true values to be equal to the predicted values. This never
happens (if it does, something has usually gone horribly wrong).

% TODO: Cut and past code

\par
The residuals of a good fit should have a symmetric distribution. A poor fit
will have a skewed distribution whilst the good fit will have a symmetric
distribution (shown below in the histograms).

% TODO: Example

% TODO: Another example

\par
What we really like though is a single value that tells us how good our model .
The most natural metric if we are using the $L2 loss$ is the root mean squared
error (RMSE).

\par
The most natural metric if we are using the $L1 loss$ is the mean absolute error
(MAE).

$$
	MAE = \overline{|\epsilon|} = \frac{1}{n}\sum|y_i - \hat{y_i}|
$$

% TODO:

\par
There are two big problems that can occur with a data set.

% TODO: Talk about drift and shift

\par
What happens if we have very little data?

$k$-fold cross validation (here $k=5$)

% TODO: Example

\par
Cross validation consists in splitting the data into $k$ subsets of equal
size to train and test on the data set $k$ times and average the results.


% TODO: Cut and paste CV scikit code

\subsubsection{Data leakage}
\par


\begin{itemize}
	\item \textbf{Target leakage:} Shouldn't have data that gives away the answer
	\item \textbf{Future leakage:} Data from the "future" that wouldn't be
	      available at the time of prediction
\end{itemize}

\par
Leakage results in overly optimistic results. In production, this can be an
issue due to not working with the same data. Future leakage is easier to solve
than target leakage, but both still have easy solutions:

\begin{itemize}
	\item No mean value imputation
	\item No scaling and splitting with future data
	      \par Cleaning  the dataframes \texttt{X\_train}, \texttt{X\_validation} and
	      \texttt{X\_test} equally, but individually
\end{itemize}

\par
To avoid future leakage via \texttt{NaN} imputation:

\begin{lstlisting}[language=Python]
# TODO
\end{lstlisting}

% Classic Andrew Ng mistake
